%% kapa: exam=1 round=33 year=2022 period=2 form=A subject=appraisal_law subject_ko=감정평가관계법규 seq=37 number=37
%% source: Q-Net 감정평가사 기출문제 (cst00309) artlSeq=5209997 file=raw/1008/5209997/2022년도 제33회 감정평가사 1차시험 2교시(감정평가 관계법규, 회계학).pdf
%% conversion: kordoc 4.24.51 -> tools/gen1.py (auto); review=proofread-sonnet; flags=none
\begin{question}{37}
도시 및 주거환경정비법령상 정비구역에서 허가를 받아야 하는 행위와 그 구체적 내용을 옳게 연결한 것은? (단, 「국토의 계획 및 이용에 관한 법률」에 따른 개발행위허가의 대상이 아닌 것을 전제로 함)
\begin{choices}
  \item 건축물의 건축 등: 「건축법」 제2조제1항제2호에 따른 건축물(가설건축물을 포함한다)의 건축, 용도변경
  \item 공작물의 설치: 농림수산물의 생산에 직접 이용되는 것으로서 국토교통부령으로 정하는 간이공작물의 설치
  \item 토석의 채취: 정비구역의 개발에 지장을 주지 아니하고 자연경관을 손상하지 아니하는 범위에서의 토석의 채취
  \item 물건을 쌓아놓는 행위: 정비구역에 존치하기로 결정된 대지에 물건을 쌓아놓는 행위
  \item 죽목의 벌채 및 식재: 관상용 죽목의 임시식재(경작지에서의 임시식재는 제외한다)
\end{choices}
\end{question}
\answer{1}
